\documentclass{article}
\usepackage{graphicx} % Required for inserting images
\usepackage{caption} % Required for \captionof
\usepackage{amsmath}
\usepackage{float}
\usepackage{comment}
\usepackage[T1]{fontenc}
\usepackage[utf8]{inputenc}
\usepackage{lmodern}
\usepackage{microtype}
\usepackage{xcolor}
\usepackage{framed}
\renewenvironment{leftbar}{%
  \def\FrameCommand{\textcolor{orange}{\vrule width 3pt} \hspace{10pt}}%
  \MakeFramed {\advance\hsize-\width \FrameRestore}}%
 {\endMakeFramed}

\title{Mechanikai mérés}
\author{Kern Luca, Vincze Csongor}
\date{2026 Április 22}

\begin{document}

\maketitle

\section*{1. Feladat: Határozza meg a spirálrugó $k^*$ direkciós nyomatékát!}

Vízmérték segítségével az asztallap vízszintesbe állítása!
Asztal sorszámának feljegyzése!

$r_{korong}= m$
$m_{edény}=4,6 g$
$\phi_{kezdeti}= °$
$m_{golyó}=4,07 g$

\begin{table}[H]
    \centering
    \begin{tabular}{|c|c|}
        \hline
        Golyók száma[db] & Szögelfordulás [°] \\ \hline
         & \\ \hline 
         & \\ \hline
         & \\ \hline 
         & \\ \hline
         & \\ \hline 
         & \\ \hline
         & \\ \hline 
         & \\ \hline
        
    \end{tabular}
    \caption{Szögelfordulás mértéke a csapágygolyók számától függően}
    \label{tab:szogelfordulas}
\end{table}


\end{document}
